\documentclass[11pt,a4paper]{article}

\usepackage[a4paper,margin=2cm,headheight=14pt]{geometry}
\usepackage{fontspec}
\IfFontExistsTF{Times New Roman}
  {\setmainfont{Times New Roman}[Ligatures=TeX]}
  {\setmainfont{TeX Gyre Termes}[Ligatures=TeX]}
\let\sffamily\rmfamily
\IfFontExistsTF{Courier New}
  {\setmonofont{Courier New}[Scale=0.9]}
  {\setmonofont{TeX Gyre Cursor}[Scale=0.95]}

\usepackage{amsmath}
\usepackage{graphicx}
\usepackage{xcolor}
\usepackage{array}
\usepackage{longtable}
\usepackage{enumitem}
\usepackage{fancyhdr}
\usepackage{titlesec}
\usepackage{float}
\usepackage{needspace}
\usepackage{tocloft}
\usepackage{etoolbox}
\usepackage{tikz}
\usetikzlibrary{arrows.meta}
\usepackage[hidelinks,unicode,pdftitle={Báo cáo Phân tích Hướng đối tượng (OOA)}]{hyperref}

\renewcommand{\contentsname}{Mục lục}
\setcounter{tocdepth}{1}
\renewcommand{\cftsecfont}{\bfseries}
\renewcommand{\cftsecleader}{\bfseries\cftdotfill{\cftdotsep}}
\renewcommand{\cftsecpagefont}{\bfseries}

\titleformat{\section}
  {\Large\bfseries}
  {\thesection.}
  {1em}
  {}
\titleformat{\subsection}
  {\large\bfseries}
  {\thesubsection}
  {1em}
  {}
\titleformat{\subsubsection}{\normalsize\bfseries\sffamily}{\thesubsubsection}{0.6em}{}
\titlespacing*{\section}{0pt}{1.6em}{0.8em}
\titlespacing*{\subsection}{0pt}{1.2em}{0.6em}
\titlespacing*{\subsubsection}{0pt}{1em}{0.5em}

\pretocmd{\subsection}{\Needspace{7\baselineskip}}{}{}

\pagestyle{fancy}
\fancyhf{}
\fancyfoot[C]{\small\thepage}
\renewcommand{\headrulewidth}{0pt}

\setlength{\parindent}{0pt}
\setlength{\parskip}{0.6em}
\setlength{\emergencystretch}{3em}
\setlist[itemize]{itemsep=3pt,topsep=4pt,parsep=0pt}

\newcommand{\code}[1]{\texttt{#1}}
\newcolumntype{L}[1]{>{\raggedright\arraybackslash}p{#1}}

\newcommand{\lthead}[1]{%
  \hline#1\\\hline\endfirsthead
  \hline#1\\\hline\endhead
}

\newenvironment{smalltable}{\begingroup\footnotesize\renewcommand{\arraystretch}{1.25}\setlength{\tabcolsep}{3pt}}{\endgroup}

\newenvironment{crc}[1]{%
  \par\vspace{0.8em}\noindent
  \begin{minipage}{\linewidth}\footnotesize
  \renewcommand{\arraystretch}{1.3}\setlength{\tabcolsep}{5pt}
  \begin{tabular}{|L{0.66\linewidth}|L{0.26\linewidth}|}
  \hline
  \multicolumn{2}{|l|}{\textbf{Lớp: \texttt{#1}}}\\
  \hline
  \textbf{Trách nhiệm (Responsibilities)} & \textbf{Cộng tác viên (Collaborators)}\\
  \hline
}{%
  \end{tabular}
  \end{minipage}\par
}

\begin{document}

\begin{titlepage}
    \begin{tikzpicture}[remember picture, overlay]
        \draw[black, line width=2pt]
            ([shift={(1cm,-1cm)}] current page.north west) rectangle
            ([shift={(-1cm,1cm)}] current page.south east);
    \end{tikzpicture}

    \setlength{\parskip}{0pt}
    \vspace*{\fill}
    \begin{center}
        \IfFileExists{logo-PTIT.jpg}{\includegraphics[width=7cm]{logo-PTIT.jpg}}{\mbox{}}\\[1.5cm]
        {\Huge \textbf{Personal Finance Management}} \\[1.5cm]
        {\huge \textbf{Nhóm 3}} \\[0.5cm]
        Lê Anh Minh -- B24DCCE180 \\
        Ứng Trọng Trình -- B24DCCE271 \\
        Nguyễn Đình Tiến -- B24DCCE264 \\
        Nguyễn Đinh Anh Quân -- B24DCCE222
    \end{center}
    \vspace*{\fill}
\end{titlepage}

\begin{center}
  {\LARGE\bfseries\sffamily Hệ thống Quản lý Tài chính Cá nhân (PFMS)\par}
  \vspace{0.6em}
\end{center}

Tài liệu này tổng hợp kết quả phân tích hướng đối tượng cho Hệ thống Quản lý Tài chính Cá nhân (PFMS), bao gồm đầy đủ phân hệ Cốt lõi (Core) và phân hệ Mở rộng (Deferred).

\tableofcontents

\pretocmd{\section}{\clearpage}{}{}

\section{Mục tiêu và phạm vi phân tích}

\subsection{Bối cảnh và mục tiêu}
Mục tiêu của tài liệu là nhận diện chính xác các thực thể đời thực trong bài toán quản lý tài chính cá nhân, xác định thông tin nội tại, ranh giới dữ liệu và trách nhiệm nghiệp vụ của từng đối tượng trước khi bước vào pha thiết kế chi tiết (OOD).

\subsection{Phân định phạm vi nghiệp vụ}
Phạm vi phân tích bao phủ toàn bộ 11 ca sử dụng trong tài liệu đặc tả yêu cầu, chia thành hai phân hệ:
\begin{itemize}
  \item Phân hệ Cốt lõi (Core): Quản lý nơi cất giữ tiền tệ (Ví tiền, Tài khoản), ghi nhận biến động dòng tiền (Giao dịch thu chi), luân chuyển tiền nội bộ bảo toàn tổng tài sản (Chuyển khoản nội bộ), và phân loại giao dịch 3 cấp (Loại luồng tiền, Nhóm nhu cầu, Danh mục chi tiết).
  \item Phân hệ Mở rộng (Deferred): Quản lý thỏa thuận vay và cho vay với lịch trả góp nhiều kỳ (Giao dịch nợ, Đối tác, Sổ nợ), và Phân bổ tích lũy tài sản (Sổ tiết kiệm, Chứng chỉ quỹ đầu tư, Quỹ mục tiêu).
\end{itemize}

\subsection{Ranh giới phân tích OOA}
Tài liệu phân định ranh giới nghiêm ngặt theo các nguyên tắc sau:
\begin{itemize}
  \item Ranh giới công nghệ và giao diện: Không đưa vào tài liệu các khái niệm về giao diện người dùng (cửa sổ JavaFX, biểu mẫu, nút bấm) hoặc cơ chế lưu trữ kỹ thuật (bảng cơ sở dữ liệu quan hệ, câu lệnh SQL, khóa ngoại cơ sở dữ liệu, tập tin lưu trữ).
  \item Ranh giới thiết kế kiến trúc: Không đưa vào các lớp kỹ thuật của pha thiết kế phần mềm như Controller, Service, Repository, DAO, hay DTO.
  \item Ranh giới điểm chạm liên phân hệ: Khi các nghiệp vụ nợ hoặc đầu tư làm thay đổi số dư ví, điểm chạm được mô tả ở mức khái quát là ví tiền nhận biến động số dư thông qua sự kiện giao dịch tài chính chung. Các lớp thuộc phân hệ cốt lõi không lưu trữ các trường dữ liệu mang tính dự phòng cho phân hệ mở rộng như mã khoản nợ, kỳ hạn sổ tiết kiệm, hay mã chứng chỉ quỹ.
\end{itemize}

\section{Cách trích xuất đối tượng từ ca sử dụng}

\subsection{Tiêu chí nhận diện và phân loại danh từ}
Toàn bộ danh từ và cụm danh từ xuất hiện trong kịch bản ca sử dụng được phân tích và phân loại thành 4 nhóm:
\begin{itemize}
  \item Đối tượng nghiệp vụ (Entity): Thực thể có định danh độc lập, nắm giữ dữ liệu và hành vi nghiệp vụ riêng, tồn tại và biến đổi theo thời gian.
  \item Đối tượng giá trị (Value Object): Kết quả tính toán tổng hợp tại một thời điểm cụ thể, không cần định danh độc lập để quản lý vòng đời.
  \item Thuộc tính của đối tượng (Attribute): Giá trị đo lường, trạng thái hoặc thông tin mô tả chi tiết trực thuộc một đối tượng khác.
  \item Yếu tố loại bỏ (Out of Scope): Thao tác giao diện, cách hiển thị đồ họa hoặc khái niệm kỹ thuật nằm ngoài bài toán phân tích nghiệp vụ.
\end{itemize}

\subsection{Bảng tổng hợp trích xuất danh từ từ kịch bản ca sử dụng}

\begin{smalltable}
\begin{longtable}{|L{0.28\linewidth}|L{0.15\linewidth}|L{0.17\linewidth}|L{0.33\linewidth}|}
\lthead{\textbf{Từ ngữ trong ca sử dụng} & \textbf{Xuất hiện ở ca sử dụng} & \textbf{Kết luận} & \textbf{Lý do}}
Ví tiền (Tài khoản, Tiền mặt, Ví nguồn, Ví đích, Ví thanh toán) & WAL-01, TXN-01, TXN-02, TRF-01, DSH-01 & Đối tượng nghiệp vụ & Nơi cất giữ tiền thực tế, có số dư biến động qua các giao dịch và chuyển khoản. Đóng vai trò nguồn, đích, thanh toán tùy ngữ cảnh. \\ \hline
Giao dịch (Khoản thu, Khoản chi) & TXN-01, TXN-02, TXN-03, CAT-01, DSH-01 & Đối tượng nghiệp vụ & Sự kiện dòng tiền đi vào hoặc đi ra khỏi ví cụ thể, có dữ liệu thời gian và số tiền cần lưu vết độc lập. \\ \hline
Chuyển khoản (Chuyển tiền nội bộ) & TRF-01 & Đối tượng nghiệp vụ & Sự kiện luân chuyển tiền tệ giữa hai ví của cùng một chủ sở hữu, bảo toàn tổng giá trị tài sản. \\ \hline
Nhóm nhu cầu (\code{Group}) & TXN-01, TXN-03, CAT-01, CAT-02, DSH-01 & Đối tượng nghiệp vụ & Cấp phân loại thứ hai, dùng để gom nhóm các danh mục chi tiết phục vụ thống kê và định hướng ngân sách. \\ \hline
Danh mục phân loại (\code{Category}, Danh mục chi tiết) & TXN-01, TXN-03, CAT-01, CAT-02 & Đối tượng nghiệp vụ & Cấp phân loại thứ ba, nhãn dán cụ thể gắn vào từng giao dịch thu chi. \\ \hline
Khoản nợ (\code{Debt}, Khoản vay, Khoản cho vay) & DBT-01 & Đối tượng nghiệp vụ & Thỏa thuận vay hoặc cho vay có định danh, có gốc và lãi suất, tồn tại và biến đổi theo lịch trả nợ. \\ \hline
Đối tác (\code{Partner}, Người vay, Người cho vay) & DBT-01, FR-DBT-03 & Đối tượng nghiệp vụ & Cá nhân hoặc tổ chức có quan hệ vay mượn với chủ tài khoản, cần định danh độc lập để tổng hợp sổ nợ theo từng người. \\ \hline
Kỳ trả nợ (\code{RepaymentInstallment}) & DBT-01 (Mở rộng theo lịch trả góp) & Đối tượng nghiệp vụ & Một kỳ thanh toán cụ thể trong lịch trả góp, quy định số tiền dự kiến và hạn chót thanh toán. \\ \hline
Giao dịch trả hoặc thu nợ (\code{DebtPayment}) & DBT-01 & Đối tượng nghiệp vụ & Sự kiện phát sinh dòng tiền thực tế để trả hoặc thu cho một kỳ nợ cụ thể, bao gồm trả nợ, thu nợ và xóa nợ. \\ \hline
Sổ nợ theo đối tác (\code{DebtLedger}, Phải thu, Phải trả) & DBT-01, FR-DBT-02, FR-DBT-03 & Đối tượng giá trị & Bảng tổng hợp dư nợ hiện tại giữa chủ tài khoản và một đối tác tại thời điểm đối soát, suy ra từ toàn bộ các khoản nợ của đối tác đó. \\ \hline
Sổ tiết kiệm (\code{Saving}) & SAV-01, GFD-01 & Đối tượng nghiệp vụ & Khoản tiền gửi tiết kiệm có kỳ hạn hoặc không kỳ hạn, sinh lãi theo thời gian và có thể liên kết vào quỹ mục tiêu. \\ \hline
Chứng chỉ quỹ (\code{FundCertificate}) & FCT-01, GFD-01 & Đối tượng nghiệp vụ & Khoản đầu tư vào chứng chỉ quỹ hoặc cổ phiếu, có giá trị biến động theo thị trường và liên kết vào quỹ mục tiêu. \\ \hline
Quỹ mục tiêu (\code{GoalFund}) & SAV-01, FCT-01, GFD-01 & Đối tượng nghiệp vụ & Mục tiêu tài chính cụ thể mà chủ sở hữu phân bổ tài sản (tiền mặt, sổ tiết kiệm, chứng chỉ quỹ) để theo dõi tiến độ hoàn thành. \\ \hline
Số dư (Ban đầu, Khả dụng, Hiện tại) & WAL-01, TXN-01, TXN-02, TRF-01, DSH-01 & Thuộc tính & Giá trị tiền tệ đo lường lượng tiền cất giữ, thuộc về cấu trúc dữ liệu của đối tượng Ví tiền. \\ \hline
Thông tin mô tả ví (Tên ví, Loại ví, Trạng thái) & WAL-01, CAT-02 & Thuộc tính & Dữ liệu cơ sở mô tả tính chất và trạng thái hoạt động của đối tượng Ví tiền. \\ \hline
Chi tiết giao dịch (Ngày giao dịch, Số tiền, Ghi chú) & TXN-01, TXN-02, TXN-03, TRF-01 & Thuộc tính & Các giá trị đo lường thời gian, tiền tệ và ghi chú cấu thành nên một Giao dịch hoặc Chuyển khoản. \\ \hline
Loại luồng tiền (\code{Type}: Thu, Chi) & TXN-01, CAT-01, CAT-02 & Thuộc tính & Thuộc tính định nghĩa phân loại gốc (Income hoặc Expense), gắn liền với Nhóm nhu cầu và Danh mục. \\ \hline
Mã đợt vay (\code{Batch}) & DBT-01 & Thuộc tính & Mã nhóm các giao dịch nợ phát sinh cùng đợt, trực thuộc đối tượng Khoản nợ. \\ \hline
Lãi suất & SAV-01, FCT-01, DBT-01 & Thuộc tính & Tỷ lệ phần trăm tính lãi trên vốn gốc, thuộc về Khoản nợ, Sổ tiết kiệm hoặc Chứng chỉ quỹ. \\ \hline
Mã định danh (\code{Id}) & Toàn bộ Use Case & Thuộc tính & Khóa định danh độc nhất của từng đối tượng thực thể trong hệ thống. \\ \hline
Chủ sở hữu tài chính (\code{Owner}, Người dùng) & Toàn bộ Use Case & Loại bỏ & Tác nhân sử dụng hệ thống. Do hệ thống thiết kế cho một người dùng cá nhân (Single-user), tác nhân không cần lưu thành đối tượng nghiệp vụ. \\ \hline
Giao diện người dùng (Màn hình, Biểu mẫu, Nút bấm, Hộp thoại) & Toàn bộ Use Case & Loại bỏ & Thuộc về tầng hiển thị giao diện, không phải là đối tượng miền nghiệp vụ. \\ \hline
Báo cáo động (Cơ cấu tài sản, Biểu đồ thu chi) & WAL-01, TXN-03, DSH-01 & Loại bỏ & Kết quả truy vấn và tổng hợp dữ liệu động từ các thực thể gốc phục vụ hiển thị màn hình, không lưu trữ độc lập. \\ \hline
Khái niệm kỹ thuật (Cơ sở dữ liệu, Bảng dữ liệu, Ghi sổ kép) & Toàn bộ Use Case & Loại bỏ & Quy tắc logic nghiệp vụ hoặc công nghệ lưu trữ, không phải là thực thể tĩnh trong mô hình phân tích. \\ \hline
\end{longtable}
\end{smalltable}

\section{Từ điển khái niệm miền}

Hệ thống tài chính cá nhân được mô hình hóa qua 13 đối tượng nghiệp vụ, gồm 12 thực thể (Entity) và 1 đối tượng giá trị (Value Object). Mọi giá trị liên quan đến tiền tệ, số dư và lãi suất bắt buộc sử dụng kiểu số thập phân chính xác (\code{BigDecimal}).

\begin{smalltable}\scriptsize\renewcommand{\arraystretch}{0.95}
\begin{longtable}{|L{0.205\linewidth}|L{0.09\linewidth}|L{0.15\linewidth}|L{0.22\linewidth}|L{0.085\linewidth}|L{0.165\linewidth}|}
\lthead{\textbf{Thuật ngữ} & \textbf{Tên tiếng Việt} & \textbf{Bản chất đời thực} & \textbf{Thông tin đối tượng nắm giữ} & \textbf{Phân loại} & \textbf{Ranh giới dữ liệu}}
\code{Wallet} & Ví tiền, Tài khoản & Nơi cất giữ tiền thực tế như tiền mặt, tài khoản ngân hàng, ví điện tử. & Mã ví, tên ví, loại ví, trạng thái hoạt động, số dư ban đầu (\code{BigDecimal}), số dư hiện tại (\code{BigDecimal}). & Entity & Quản lý số dư tự thân. Nhận biến động tiền tệ qua sự kiện giao dịch, không lưu thông tin chi tiết về nợ hay khoản đầu tư. \\ \hline
\code{Transaction} & Giao dịch thu chi & Sự kiện tiền tệ biến động làm tăng hoặc giảm số dư của một ví cụ thể. & Mã giao dịch, ngày giao dịch, số tiền (\code{BigDecimal}), loại luồng tiền (Income hoặc Expense), ghi chú, mã ví, mã danh mục. & Entity & Chỉ lưu vết lịch sử giao dịch. Không tự động sửa số dư ví mà cung cấp thông tin để Ví cập nhật. \\ \hline
\code{Transfer} & Chuyển khoản nội bộ & Sự kiện luân chuyển tiền tệ giữa hai ví của cùng một chủ tài khoản, bảo toàn tổng tài sản. & Mã chuyển khoản, ngày chuyển, số tiền (\code{BigDecimal}), mức phí chuyển (\code{BigDecimal}), ghi chú, mã ví nguồn, mã ví đích. & Entity & Chỉ theo dõi luồng chuyển tiền giữa hai ví. Phí phát sinh (nếu có) được tách thành một giao dịch chi tiêu riêng. \\ \hline
\code{Group} & Nhóm nhu cầu & Cấp phân loại thu chi nhóm 2 (Nhu cầu thiết yếu, Đầu tư, Tiết kiệm, Hưởng thụ). & Mã nhóm, tên nhóm, loại luồng tiền gốc (Income hoặc Expense). & Entity & Quản lý cây phân loại cấp 2, không lưu trữ tổng tiền giao dịch. \\ \hline
\code{Category} & Danh mục phân loại & Nhãn phân loại chi tiết cấp 3 gắn trực tiếp vào từng giao dịch thu chi. & Mã danh mục, tên danh mục, biểu tượng, màu sắc, trạng thái hiển thị, mã nhóm nhu cầu cha. & Entity & Quản lý nhãn phân loại phục vụ thống kê, hoàn toàn độc lập với dữ liệu số dư tiền tệ. \\ \hline
\code{Debt} & Khoản nợ & Thỏa thuận vay hoặc cho vay cụ thể với một đối tác, có gốc và lãi suất cố định. & Mã khoản nợ, loại nghiệp vụ (Lend hoặc Borrow), mã đối tác, số tiền gốc (\code{BigDecimal}), lãi suất năm (\code{BigDecimal}), ngày tạo, tổng số kỳ trả góp, trạng thái khoản nợ. & Entity & Quản lý thông tin thỏa thuận nợ và các kỳ trả góp trực thuộc, không lưu trữ mã ví thực hiện thanh toán. \\ \hline
\code{RepaymentInstallment} & Kỳ trả nợ & Một kỳ hạn cụ thể trong lịch trả góp của khoản nợ. & Số thứ tự kỳ, hạn thanh toán, số tiền dự kiến phải trả gồm gốc và lãi đơn (\code{BigDecimal}), trạng thái thanh toán của kỳ. & Entity & Theo dõi tiến độ thanh toán của từng kỳ, không trực tiếp quản lý dòng tiền tổng của toàn bộ khoản nợ. \\ \hline
\code{DebtPayment} & Giao dịch trả hoặc thu nợ & Sự kiện thanh toán tiền thực tế phát sinh cho một kỳ trả nợ cụ thể. & Mã giao dịch nợ, ngày thanh toán, số tiền thực trả (\code{BigDecimal}), loại thao tác (Collect, Repay, Write-off), mã kỳ trả nợ, ghi chú. & Entity & Lưu vết sự kiện thanh toán nợ, là điểm chạm kích hoạt thay đổi số dư ví tương ứng. \\ \hline
\code{Partner} & Đối tác công nợ & Cá nhân hoặc tổ chức có quan hệ vay mượn với chủ tài khoản. & Mã đối tác, tên đối tác, loại mối quan hệ (Gia đình, Bạn bè, Công việc, Khác), thông tin liên hệ. & Entity & Quản lý danh mục đối tác, không lưu trữ số dư nợ trực tiếp trên bản ghi đối tác. \\ \hline
\code{DebtLedger} & Sổ nợ theo đối tác & Bảng tổng hợp số dư công nợ giữa chủ tài khoản và một đối tác tại thời điểm đối soát. & Đối tác, tổng số tiền phải thu (\code{BigDecimal}), tổng số tiền phải trả (\code{BigDecimal}), dư nợ ròng (\code{BigDecimal}). & Value Object & Tính toán động từ toàn bộ các khoản nợ của đối tác, không tự lưu trữ trạng thái độc lập. \\ \hline
\code{Saving} & Sổ tiết kiệm & Khoản tiền gửi tiết kiệm tại tổ chức tín dụng có hoặc không có kỳ hạn, sinh lãi theo thời gian. & Mã sổ, mã quỹ mục tiêu liên kết, số tiền gửi ban đầu (\code{BigDecimal}), kỳ hạn gửi, lãi suất năm (\code{BigDecimal}), ngày gửi, trạng thái sổ. & Entity & Quản lý thông tin khoản tiền gửi tiết kiệm tự thân, không lưu thông tin chi tiết của ví nguồn. \\ \hline
\code{FundCertificate} & Chứng chỉ quỹ đầu tư & Khoản đầu tư vào chứng chỉ quỹ mở hoặc cổ phiếu, có giá trị thị trường biến động liên tục. & Mã chứng chỉ, mã quỹ mục tiêu liên kết, mã giao dịch chứng chỉ, loại tài sản (Cổ phiếu, Trái phiếu, Cân bằng), vốn gốc (\code{BigDecimal}), ngày mua, giá trị thị trường hiện tại (\code{BigDecimal}). & Entity & Quản lý danh mục đầu tư và giá trị thị trường, không lưu lịch sử giao dịch của ví thanh toán. \\ \hline
\code{GoalFund} & Quỹ mục tiêu & Quỹ tài chính do chủ tài khoản thiết lập để tích lũy tài sản cho mục tiêu cụ thể. & Mã quỹ, tên quỹ, số tiền mục tiêu (\code{BigDecimal}), số tiền mặt phân bổ trực tiếp (\code{BigDecimal}), ngày bắt đầu, ngày dự kiến hoàn thành. & Entity & Tổng hợp giá trị tài sản từ tiền mặt phân bổ, các sổ tiết kiệm và các chứng chỉ quỹ liên kết, không lưu trùng lặp số dư của tài sản thành phần. \\ \hline
\end{longtable}
\end{smalltable}

\section{Mối quan hệ nghiệp vụ giữa các đối tượng}

\subsection{Bảng tổng hợp quan hệ nghiệp vụ}

\begin{smalltable}\scriptsize
\begin{longtable}{|L{0.205\linewidth}|L{0.09\linewidth}|L{0.205\linewidth}|L{0.205\linewidth}|L{0.205\linewidth}|}
\lthead{\textbf{Đối tượng A} & \textbf{Quan hệ} & \textbf{Đối tượng B} & \textbf{Số lượng (Bản số)} & \textbf{Ý nghĩa nghiệp vụ}}
\code{Wallet} & chứa & \code{Transaction} & 1 \code{Wallet} chứa 0 hoặc nhiều \code{Transaction} & Mỗi khoản thu hoặc chi bắt buộc phải gắn với một ví cụ thể để xác định nguồn tiền biến động. \\ \hline
\code{Transaction} & thuộc & \code{Category} & 1 \code{Transaction} thuộc đúng 1 \code{Category} & Mọi giao dịch thu chi đều cần nhãn phân loại phục vụ thống kê chi tiêu và lập ngân sách. \\ \hline
\code{Category} & thuộc & \code{Group} & 1 \code{Category} thuộc đúng 1 \code{Group} & Các danh mục chi tiết được gom nhóm vào nhóm nhu cầu lớn tương ứng với luồng tiền gốc. \\ \hline
\code{Transfer} & luân chuyển giữa & \code{Wallet} & 1 \code{Transfer} liên kết đúng 2 \code{Wallet} & Chuyển tiền nội bộ làm giảm số dư ví nguồn và tăng số dư ví đích của cùng một người dùng. \\ \hline
\code{Transfer} & phát sinh phí qua & \code{Transaction} & 1 \code{Transfer} kích hoạt 0 hoặc 1 \code{Transaction} & Trường hợp chuyển khoản có tính phí dịch vụ, phí chuyển được ghi nhận thành một giao dịch chi độc lập trừ vào ví nguồn. \\ \hline
\code{Partner} & tham gia & \code{Debt} & 1 \code{Partner} có 0 hoặc nhiều \code{Debt} & Một đối tác có thể có nhiều thỏa thuận vay hoặc cho vay với chủ tài khoản theo thời gian. \\ \hline
\code{Debt} & chia thành & \code{RepaymentInstallment} & 1 \code{Debt} có 1 hoặc nhiều \code{RepaymentInstallment} & Khoản nợ được chia thành nhiều kỳ thanh toán theo lịch trả góp đã thỏa thuận. \\ \hline
\code{RepaymentInstallment} & nhận thanh toán từ & \code{DebtPayment} & 1 \code{RepaymentInstallment} có 0 hoặc nhiều \code{DebtPayment} & Một kỳ hạn có thể được thanh toán một lần, nhiều lần, hoặc xử lý xóa nợ khi không đòi được. \\ \hline
\code{Partner} & tổng hợp đối soát & \code{DebtLedger} & 1 \code{Partner} tương ứng đúng 1 \code{DebtLedger} & Sổ nợ được tính toán tức thời từ toàn bộ các khoản nợ của đối tác đó để xác định dư nợ ròng. \\ \hline
\code{GoalFund} & phân bổ vào & \code{Saving} & 1 \code{GoalFund} chứa 0 hoặc nhiều \code{Saving} & Quỹ mục tiêu gom nhóm các sổ tiết kiệm phục vụ cùng một mục tiêu tài chính. \\ \hline
\code{GoalFund} & phân bổ vào & \code{FundCertificate} & 1 \code{GoalFund} chứa 0 hoặc nhiều \code{FundCertificate} & Quỹ mục tiêu tích hợp các khoản đầu tư chứng chỉ quỹ để theo dõi tiến độ tổng thể. \\ \hline
\code{DebtPayment} & tác động số dư & \code{Wallet} & Nhiều \code{DebtPayment} tác động đến 1 \code{Wallet} & Điểm chạm liên phân hệ: Mỗi giao dịch trả nợ hoặc thu nợ làm biến động số dư của chiếc ví được chọn thanh toán. \\ \hline
\end{longtable}
\end{smalltable}

\subsection{Điểm chạm liên phân hệ (Inter-subsystem Touchpoints)}
Các phân hệ trong hệ thống tương tác với nhau qua ba điểm chạm dữ liệu chính:
\begin{itemize}
  \item Điểm chạm thanh toán nợ: Khi phát sinh một sự kiện trả nợ (giảm tiền) hoặc thu nợ (tăng tiền) trong \code{DebtPayment}, số dư của \code{Wallet} được điều chỉnh thông qua một biến động dòng tiền tương ứng.
  \item Điểm chạm phân bổ quỹ mục tiêu: Khi chủ tài khoản trích tiền mặt từ \code{Wallet} để đưa vào \code{GoalFund} (khoản mục \code{funds\_cash}), số dư ví thanh toán giảm tương ứng với số tiền phân bổ.
  \item Điểm chạm gửi tiết kiệm và mua chứng chỉ quỹ: Khi khởi tạo một \code{Saving} hoặc \code{FundCertificate}, dòng tiền xuất phát từ một \code{Wallet} xác định. Khi sổ tiết kiệm đáo hạn hoặc chứng chỉ quỹ được bán, dòng tiền quay trở lại ví đích được chỉ định.
\end{itemize}

\subsection{Mô hình trực quan hai tầng (Tiered Visual Modeling)}

Nhằm đảm bảo tính trực quan và tuân thủ giới hạn nhận thức tối đa 9 node trong một biểu đồ, kiến trúc quan hệ được phân rã thành hai tầng mô hình hóa:

\Needspace{9cm}
\subsubsection*{Sơ đồ Tầng 1: Tổng quan tương tác giữa các phân hệ nghiệp vụ}

Biểu đồ mô tả luồng tương tác cấp cao giữa nơi giữ tiền trung tâm và các phân hệ chức năng:

\begin{center}
\begin{tikzpicture}[
  box/.style={draw, rounded corners=2pt, thick, minimum width=3.9cm, minimum height=1.3cm,
              align=center, font=\footnotesize, text width=3.6cm},
  lbl/.style={font=\scriptsize, fill=white, inner sep=1.5pt, align=center, text width=3.0cm},
  sflow/.style={-{Stealth[length=2.5mm]}, thick},
  dflow/.style={-{Stealth[length=2.5mm]}, thick, dash pattern=on 4pt off 3pt}
]
  \node[box] (wal)  at (0,0)      {\textbf{Wallet}\\Ví tiền \& Tài khoản};
  \node[box] (core) at (-4.3,3.8) {\textbf{Core}\\Giao dịch thu chi \& Chuyển khoản};
  \node[box] (debt) at (4.3,3.8)  {\textbf{Debt}\\Công nợ \& Lịch trả góp};
  \node[box] (fund) at (-4.3,-3.8){\textbf{Funds}\\Sổ tiết kiệm \& Chứng chỉ quỹ};
  \node[box] (goal) at (4.3,-3.8) {\textbf{GoalFund}\\Quỹ mục tiêu tích lũy};

  \draw[sflow] (core) -- node[lbl, midway] {Cập nhật thu/chi và chuyển khoản} (wal);
  \draw[dflow] (debt) -- node[lbl, midway] {DebtPayment làm biến động số dư} (wal);
  \draw[dflow] (fund) -- node[lbl, midway] {Nguồn tiền gửi và tiền rút} (wal);
  \draw[dflow] (goal) -- node[lbl, midway] {Phân bổ tiền mặt trực tiếp} (wal);
  \draw[sflow] (fund) -- node[lbl, midway] {Tích hợp giá trị tài sản} (goal);
\end{tikzpicture}
\end{center}

\Needspace{9cm}
\subsubsection*{Sơ đồ Tầng 2a: Chi tiết cấu trúc phân hệ Cốt lõi}

Biểu đồ mô tả cấu trúc quan hệ nội tại của 5 đối tượng cốt lõi:

\begin{center}
\begin{tikzpicture}[
  box/.style={draw, rounded corners=2pt, thick, minimum width=3.4cm, minimum height=1.1cm,
              align=center, font=\footnotesize},
  lbl/.style={font=\scriptsize, fill=white, inner sep=1.5pt, align=center, text width=2.4cm},
  sflow/.style={-{Stealth[length=2.5mm]}, thick},
  dflow/.style={-{Stealth[length=2.5mm]}, thick, dash pattern=on 4pt off 3pt}
]
  \node[box] (grp) at (0,6.4)   {\textbf{Group}\\Nhóm nhu cầu};
  \node[box] (cat) at (0,3.2)   {\textbf{Category}\\Danh mục};
  \node[box] (txn) at (0,0)     {\textbf{Transaction}\\Giao dịch};
  \node[box] (wal) at (6.8,3.2) {\textbf{Wallet}\\Ví tiền};
  \node[box] (trf) at (6.8,0)   {\textbf{Transfer}\\Chuyển khoản};

  \draw[sflow] (grp) -- node[lbl, midway] {1 chứa nhiều} (cat);
  \draw[sflow] (cat) -- node[lbl, midway] {1 dán nhãn nhiều} (txn);
  \draw[sflow] (wal) -- node[lbl, midway] {1 chứa nhiều} (txn);
  \draw[sflow] (wal) -- node[lbl, midway] {2 tham gia 1} (trf);
  \draw[dflow] (trf) -- node[lbl, midway] {Phát sinh phí dịch vụ} (txn);
\end{tikzpicture}
\end{center}

\Needspace{9cm}
\subsubsection*{Sơ đồ Tầng 2b: Chi tiết cấu trúc phân hệ Mở rộng}

Biểu đồ mô tả cấu trúc quan hệ của 8 đối tượng thuộc phân hệ Công nợ và Quỹ tài chính:

\begin{center}
\begin{tikzpicture}[
  box/.style={draw, rounded corners=2pt, thick, minimum height=1.0cm,
              align=center, font=\footnotesize\bfseries},
  lbl/.style={font=\scriptsize, fill=white, inner sep=1.5pt, align=center, text width=1.8cm},
  sflow/.style={-{Stealth[length=2.5mm]}, thick}
]
  \node[box, minimum width=2.4cm] (par) at (0,0)     {\code{Partner}};
  \node[box, minimum width=2.4cm] (dbt) at (4.9,0)   {\code{Debt}};
  \node[box, minimum width=4.0cm] (ins) at (10.3,0)  {\code{RepaymentInstallment}};
  \node[box, minimum width=3.0cm] (pay) at (10.3,-2.8) {\code{DebtPayment}};
  \node[box, minimum width=2.4cm] (led) at (0,-2.8)  {\code{DebtLedger}};

  \node[box, minimum width=2.4cm] (gf)  at (0,-6.2)   {\code{GoalFund}};
  \node[box, minimum width=3.0cm] (sav) at (6.0,-5.2) {\code{Saving}};
  \node[box, minimum width=3.0cm] (fc)  at (6.0,-7.2) {\code{FundCertificate}};

  \draw[sflow] (par) -- node[lbl, midway] {1 có nhiều} (dbt);
  \draw[sflow] (dbt) -- node[lbl, midway] {1 chia thành nhiều} (ins);
  \draw[sflow] (ins) -- node[lbl, midway] {1 nhận nhiều} (pay);
  \draw[sflow] (par) -- node[lbl, midway] {1 tổng hợp thành} (led);
  \draw[sflow] (gf)  -- node[lbl, text width=2.6cm, midway] {1 phân bổ nhiều} (sav);
  \draw[sflow] (gf)  -- node[lbl, text width=2.6cm, midway] {1 tích hợp nhiều} (fc);
\end{tikzpicture}
\end{center}

\section{Phân bổ trách nhiệm bằng thẻ CRC}

Phương pháp thẻ CRC (Class, Responsibility, Collaborator) được sử dụng để phân định rõ trách nhiệm nội tại của từng đối tượng. Cột trách nhiệm phân biệt giữa việc đối tượng biết (Knowing) và việc đối tượng tự thực hiện (Doing). Cột cộng tác viên xác định các đối tượng khác cần phối hợp.

\subsection{Thẻ CRC phân hệ Cốt lõi}

\begin{crc}{Wallet}
Biết mã định danh, tên ví, loại ví và trạng thái hoạt động. & \\ \hline
Biết số dư ban đầu và duy trì số dư khả dụng hiện tại. & \\ \hline
Kiểm tra số dư khả dụng trước khi chi tiêu hoặc chuyển tiền. & \\ \hline
Tự động cập nhật số dư khi nhận biến động giao dịch hoặc chuyển khoản. & \code{Transaction}, \code{Transfer} \\ \hline
Ngăn chặn hành vi đóng hoặc xóa ví khi số dư hiện tại khác 0. & \\ \hline
\end{crc}

\begin{crc}{Transaction}
Biết mã giao dịch, ngày thực hiện, số tiền, loại luồng tiền và ghi chú. & \\ \hline
Xác thực số tiền giao dịch phải là số dương lớn hơn 0. & \\ \hline
Liên kết với danh mục phân loại chi tiết tương ứng. & \code{Category} \\ \hline
Cung cấp dữ liệu biến động dòng tiền để ví tự cập nhật số dư. & \code{Wallet} \\ \hline
\end{crc}

\begin{crc}{Transfer}
Biết mã chuyển khoản, ngày chuyển, số tiền, mức phí dịch vụ và ghi chú. & \\ \hline
Điều phối biến động giảm số dư ở ví nguồn và tăng số dư ở ví đích. & \code{Wallet} \\ \hline
Kích hoạt tạo giao dịch chi tiêu độc lập khi có phát sinh phí chuyển tiền. & \code{Transaction} \\ \hline
Xác thực ví nguồn và ví đích không được trùng nhau. & \code{Wallet} \\ \hline
\end{crc}

\begin{crc}{Group}
Biết mã nhóm, tên nhóm và loại luồng tiền gốc (Income hoặc Expense). & \\ \hline
Quản lý danh sách các danh mục chi tiết trực thuộc nhóm. & \code{Category} \\ \hline
Kiểm tra tính nhất quán về loại luồng tiền của các danh mục con. & \code{Category} \\ \hline
\end{crc}

\begin{crc}{Category}
Biết mã danh mục, tên danh mục, biểu tượng, màu sắc và trạng thái hiển thị. & \\ \hline
Cung cấp nhãn phân loại phục vụ thống kê báo cáo cho giao dịch. & \code{Transaction} \\ \hline
Tham chiếu đến nhóm nhu cầu cấp trên mà danh mục trực thuộc. & \code{Group} \\ \hline
\end{crc}

\subsection{Thẻ CRC phân hệ Mở rộng}

\begin{crc}{Debt}
Biết mã khoản nợ, đối tác, số tiền vốn gốc, lãi suất năm và loại nghiệp vụ. & \code{Partner} \\ \hline
Tự sinh lịch trả góp gồm các kỳ trả nợ khi khoản nợ được khởi tạo. & \code{RepaymentInstallment} \\ \hline
Tự cập nhật trạng thái khoản nợ khi các kỳ hạn hoàn thành thanh toán. & \code{RepaymentInstallment} \\ \hline
Xác thực số tiền vốn gốc lớn hơn 0 và lãi suất không âm. & \\ \hline
\end{crc}

\begin{crc}{RepaymentInstallment}
Biết số thứ tự kỳ, hạn thanh toán và số tiền dự kiến gồm gốc và lãi. & \\ \hline
Tự xác định trạng thái theo thời gian hiện tại (Chưa đến hạn hoặc Quá hạn). & \\ \hline
Tự cập nhật trạng thái Đã trả khi nhận đủ số tiền từ các lần thanh toán. & \code{DebtPayment} \\ \hline
\end{crc}

\begin{crc}{DebtPayment}
Biết ngày thanh toán, số tiền thực trả, thao tác xử lý và kỳ hạn liên quan. & \code{RepaymentInstallment} \\ \hline
Xác thực số tiền thanh toán thực tế phải lớn hơn 0. & \\ \hline
Kích hoạt biến động số dư trên chiếc ví được chỉ định thanh toán. & \code{Wallet} \\ \hline
\end{crc}

\begin{crc}{Partner}
Biết mã đối tác, tên đối tác, loại mối quan hệ và thông tin liên hệ. & \\ \hline
Cung cấp danh sách toàn bộ các khoản vay và cho vay liên quan đến mình. & \code{Debt} \\ \hline
\end{crc}

\begin{crc}{DebtLedger}
Biết đối tác cần đối soát công nợ. & \code{Partner} \\ \hline
Tính tổng số tiền phải thu từ các khoản cho vay trừ đi số tiền đã thu hồi. & \code{Debt} \\ \hline
Tính tổng số tiền phải trả từ các khoản đi vay trừ đi số tiền đã trả. & \code{Debt} \\ \hline
Tính số dư nợ ròng giữa chủ tài khoản và đối tác tại thời điểm hiện tại. & \code{Debt} \\ \hline
\end{crc}

\begin{crc}{Saving}
Biết mã sổ, vốn gốc, kỳ hạn gửi, lãi suất năm và ngày gửi. & \\ \hline
Tự tính số dư hiện tại gồm tiền gốc và lãi tích lũy theo thời gian nắm giữ. & \\ \hline
Tự tính tiến độ hoàn thành kỳ hạn gửi tiết kiệm. & \\ \hline
\end{crc}

\begin{crc}{FundCertificate}
Biết mã chứng chỉ, loại tài sản đầu tư, vốn gốc, ngày mua và giá thị trường. & \\ \hline
Tự tính lãi hoặc lỗ bằng hiệu số giữa giá trị thị trường hiện tại và vốn gốc. & \\ \hline
Tự tính tỷ suất sinh lời năm hóa dựa trên thời gian nắm giữ. & \\ \hline
\end{crc}

\begin{crc}{GoalFund}
Biết mã quỹ, tên quỹ mục tiêu, số tiền mục tiêu và tiền mặt phân bổ. & \\ \hline
Tự tính tổng giá trị quỹ bằng tổng tiền mặt, các sổ tiết kiệm và chứng chỉ quỹ. & \code{Saving}, \code{FundCertificate} \\ \hline
Tự tính tỷ lệ phần trăm tiến độ đạt được so với số tiền mục tiêu đề ra. & \\ \hline
\end{crc}

\section{Quy tắc nghiệp vụ và ràng buộc dữ liệu}

\subsection{Quy tắc toàn hệ thống và tính toán chính xác}

\begin{smalltable}
\begin{longtable}{|L{0.12\linewidth}|L{0.17\linewidth}|L{0.22\linewidth}|L{0.42\linewidth}|}
\lthead{\textbf{Mã quy tắc} & \textbf{Tên quy tắc} & \textbf{Đối tượng áp dụng} & \textbf{Nội dung quy tắc}}
\code{BR-SYS-01} & Độ chính xác tiền tệ & Toàn hệ thống & Mọi giá trị tiền tệ, số dư tài khoản và lãi suất bắt buộc sử dụng kiểu số thập phân chính xác (\code{BigDecimal}). Kết quả tính toán hiển thị được làm tròn đến 2 chữ số thập phân theo quy chuẩn kế toán. \\ \hline
\end{longtable}
\end{smalltable}

\subsection{Quy tắc phân hệ Ví tiền và Luân chuyển dòng tiền}

\begin{smalltable}
\begin{longtable}{|L{0.12\linewidth}|L{0.17\linewidth}|L{0.22\linewidth}|L{0.42\linewidth}|}
\lthead{\textbf{Mã quy tắc} & \textbf{Tên quy tắc} & \textbf{Đối tượng áp dụng} & \textbf{Nội dung quy tắc}}
\code{BR-WAL-01} & Cân bằng số dư ví & \code{Wallet} & Số dư hiện tại = Số dư ban đầu + Tổng Thu $-$ Tổng Chi + Tổng chuyển đến $-$ Tổng chuyển đi $\pm$ Biến động từ phân hệ nợ và quỹ đầu tư. \\ \hline
\code{BR-WAL-02} & Ràng buộc số dư không âm & \code{Wallet} & Ví tiền mặt và tài khoản ngân hàng thông thường không được phép có số dư nhỏ hơn 0. Ví tín dụng được phép có số dư âm trong phạm vi hạn mức thấu chi được cấp. \\ \hline
\code{BR-WAL-03} & An toàn khi đóng ví & \code{Wallet} & Hệ thống ngăn chặn xóa hoặc lưu trữ ví nếu số dư hiện tại khác 0. Chủ sở hữu phải chuyển hết tiền sang ví khác trước khi đóng ví. \\ \hline
\code{BR-WAL-04} & Kiểm tra khả năng thanh toán & \code{Wallet}, \code{Transfer} & Hệ thống bắt buộc kiểm tra số dư khả dụng của ví nguồn trước khi xác nhận giao dịch chi tiêu hoặc chuyển tiền (ngoại trừ ví tín dụng). \\ \hline
\code{BR-TXN-01} & Số tiền giao dịch dương & \code{Transaction} & Giá trị số tiền ghi nhận trong mỗi giao dịch thu chi bắt buộc phải là số dương lớn hơn 0. Bản chất làm tăng hoặc giảm số dư phụ thuộc vào Loại giao dịch (Income hoặc Expense). \\ \hline
\code{BR-TXN-02} & Nhất quán loại luồng tiền & \code{Transaction} & Mỗi giao dịch bắt buộc phải gắn với danh mục thuộc cùng nhóm luồng tiền (giao dịch Thu chỉ chọn danh mục Thu, giao dịch Chi chỉ chọn danh mục Chi). \\ \hline
\code{BR-CAT-01} & Toàn vẹn lịch sử danh mục & \code{Category} & Không cho phép xóa vĩnh viễn danh mục đã có giao dịch tham chiếu trong lịch sử. Hệ thống áp dụng cơ chế vô hiệu hóa (soft-delete) để ẩn danh mục khỏi danh sách lựa chọn mới. \\ \hline
\code{BR-CAT-02} & Tính khả dụng của danh mục & \code{Category}, \code{Transaction} & Giao dịch mới chỉ được phép dán nhãn bằng các danh mục đang ở trạng thái hoạt động (Active). \\ \hline
\code{BR-TRF-01} & Nguyên tắc bảo toàn chuyển khoản & \code{Transfer} & Số tiền trừ ở ví nguồn phải bằng chính xác số tiền cộng vào ví đích. Ví nguồn và ví đích bắt buộc không được trùng nhau. \\ \hline
\code{BR-TRF-02} & Tách biệt phí chuyển khoản & \code{Transfer} & Phí dịch vụ phát sinh trong giao dịch chuyển khoản được ghi nhận thành một giao dịch chi tiêu độc lập trừ trực tiếp vào ví nguồn. \\ \hline
\code{BR-TRF-03} & Đồng nhất loại tiền tệ & \code{Transfer} & Ví nguồn và ví đích phải sử dụng cùng một loại đơn vị tiền tệ. Trường hợp khác loại tiền tệ, hệ thống yêu cầu xác định tỷ giá quy đổi trước khi thực hiện. \\ \hline
\end{longtable}
\end{smalltable}

\subsection{Quy tắc phân hệ Công nợ và Quỹ tài chính}

\begin{smalltable}
\begin{longtable}{|L{0.12\linewidth}|L{0.17\linewidth}|L{0.22\linewidth}|L{0.42\linewidth}|}
\lthead{\textbf{Mã quy tắc} & \textbf{Tên quy tắc} & \textbf{Đối tượng áp dụng} & \textbf{Nội dung quy tắc}}
\code{BR-DBT-01} & Vốn gốc khoản nợ dương & \code{Debt} & Số tiền vốn gốc của thỏa thuận vay hoặc cho vay bắt buộc phải lớn hơn 0. \\ \hline
\code{BR-DBT-02} & Lãi đơn cố định theo ngày & \code{Debt} & Tiền lãi mỗi kỳ được tính theo công thức lãi đơn: $\text{Tiền lãi} = \text{Gốc} \times \text{Lãi suất năm} \times \dfrac{\text{Số ngày trong kỳ}}{365}$. Hệ thống không áp dụng lãi kép cho các khoản vay cá nhân. \\ \hline
\code{BR-DBT-03} & Khớp nối tổng kỳ trả nợ & \code{RepaymentInstallment} & Tổng số tiền dự kiến của toàn bộ các kỳ trong lịch trả nợ phải bằng đúng số tiền vốn gốc cộng tổng tiền lãi dự kiến của khoản nợ. \\ \hline
\code{BR-DBT-04} & Giới hạn thanh toán kỳ & \code{DebtPayment} & Tổng số tiền các lần thanh toán thực tế cho một kỳ hạn không được vượt quá số tiền dự kiến phải trả của kỳ đó. \\ \hline
\code{BR-DBT-05} & Tính chất trạng thái kết thúc & \code{Debt} & Khoản nợ đã chuyển sang trạng thái Đã tất toán hoặc Đã xóa nợ sẽ bị đóng băng, không được tạo thêm kỳ trả góp hoặc giao dịch thanh toán mới. \\ \hline
\code{BR-PRT-01} & Toàn vẹn đối tác công nợ & \code{Partner} & Không cho phép xóa đối tác khỏi hệ thống nếu vẫn còn khoản nợ chưa tất toán gắn liền với đối tác đó. \\ \hline
\code{BR-SAV-01} & Toàn vẹn kỳ hạn tiết kiệm & \code{Saving} & Giá trị kỳ hạn gửi tiết kiệm bắt buộc phải nằm trong danh mục các kỳ hạn đã được cấu hình sẵn trong hệ thống. \\ \hline
\code{BR-FCT-01} & Vốn gốc đầu tư hợp lệ & \code{FundCertificate} & Vốn gốc đầu tư chứng chỉ quỹ bắt buộc phải lớn hơn 0, và ngày mua không được nằm ở tương lai so với ngày hiện tại. \\ \hline
\code{BR-GFD-01} & An toàn tính toán tiến độ quỹ & \code{GoalFund} & Trường hợp số tiền mục tiêu của quỹ bằng 0, tỷ lệ phần trăm tiến độ được hiển thị ở trạng thái chưa xác định (N/A) thay vì thực hiện phép chia cho 0. \\ \hline
\code{BR-GFD-02} & Toàn vẹn tài sản quỹ mục tiêu & \code{GoalFund} & Mỗi sổ tiết kiệm hoặc chứng chỉ quỹ chỉ được phép liên kết với tối đa một quỹ mục tiêu duy nhất tại một thời điểm. \\ \hline
\end{longtable}
\end{smalltable}

\section{Vòng đời nghiệp vụ của các đối tượng chính}

\subsection{Phân định tính chất vòng đời đối tượng}
Trong hệ thống tài chính cá nhân, các đối tượng được phân loại thành hai nhóm hành vi:
\begin{itemize}
  \item Nhóm sự kiện bất biến và danh mục tĩnh: Các đối tượng \code{Transaction} và \code{Transfer} là các sự kiện tài chính đã phát sinh, một khi đã lưu vết thì mang tính bất biến và không có chuyển đổi trạng thái rời rạc. Đối tượng \code{Wallet} duy trì số dư khả dụng liên tục theo thời gian thực mà không qua các bước trạng thái. Các đối tượng \code{Partner}, \code{DebtLedger}, \code{FundCertificate}, và \code{GoalFund} là danh mục tĩnh hoặc giá trị tài sản biến động liên tục theo thị trường. Do đó, các đối tượng này áp dụng ghi chú Không áp dụng (N/A) đối với máy trạng thái rời rạc theo quy định tại khung phân tích OOA.
  \item Nhóm thực thể có vòng đời trạng thái rời rạc: Gồm \code{Debt} (Khoản nợ), \code{RepaymentInstallment} (Kỳ trả nợ), và \code{Saving} (Sổ tiết kiệm). Các thực thể này trải qua nhiều giai đoạn nghiệp vụ rõ ràng từ khi khởi tạo, thực hiện nghĩa vụ cho đến khi kết thúc.
\end{itemize}

\subsection{Vòng đời Khoản nợ (\texttt{Debt})}

\begin{smalltable}
\begin{longtable}{|L{0.15\linewidth}|L{0.19\linewidth}|L{0.19\linewidth}|L{0.13\linewidth}|L{0.27\linewidth}|}
\lthead{\textbf{Trạng thái hiện tại} & \textbf{Sự kiện kích hoạt} & \textbf{Điều kiện chuyển trạng thái} & \textbf{Trạng thái tiếp theo} & \textbf{Ý nghĩa nghiệp vụ}}
Mới phát sinh & Hệ thống sinh lịch trả góp & Khoản nợ vừa được tạo lập thành công & Đang trả dần & Bắt đầu theo dõi nghĩa vụ trả nợ theo từng kỳ hạn. \\ \hline
Đang trả dần & Ghi nhận một giao dịch thanh toán nợ (\code{DebtPayment}) & Vẫn còn kỳ trả nợ chưa thanh toán đủ & Đang trả dần & Số dư nợ gốc và lãi giảm dần nhưng nghĩa vụ nợ chưa hoàn tất. \\ \hline
Đang trả dần & Kỳ hạn cuối cùng được thanh toán đủ & Tổng số tiền đã thanh toán bằng gốc cộng tổng lãi dự kiến & Đã tất toán & Khoản nợ hoàn tất trọn vẹn, chấm dứt toàn bộ nghĩa vụ nợ. \\ \hline
Đang trả dần hoặc Mới phát sinh & Chủ tài khoản xác nhận xử lý xóa nợ & Không có khả năng thu hồi hoặc có thỏa thuận miễn giảm nợ & Đã xóa nợ & Kết thúc theo dõi khoản nợ mà không thu hoặc trả đủ toàn bộ số tiền. \\ \hline
\end{longtable}
\end{smalltable}

\subsection{Vòng đời Kỳ trả nợ (\texttt{RepaymentInstallment})}

\begin{smalltable}
\begin{longtable}{|L{0.15\linewidth}|L{0.19\linewidth}|L{0.19\linewidth}|L{0.13\linewidth}|L{0.27\linewidth}|}
\lthead{\textbf{Trạng thái hiện tại} & \textbf{Sự kiện kích hoạt} & \textbf{Điều kiện chuyển trạng thái} & \textbf{Trạng thái tiếp theo} & \textbf{Ý nghĩa nghiệp vụ}}
Chưa đến hạn & Đến hạn thanh toán mà chưa nhận đủ tiền & Ngày hiện tại lớn hơn hoặc bằng ngày đến hạn của kỳ & Quá hạn & Cảnh báo nghĩa vụ thanh toán đã bị trễ hạn cần xử lý. \\ \hline
Chưa đến hạn hoặc Quá hạn & Nhận đủ tiền thanh toán qua các giao dịch nợ & Tổng số tiền thanh toán thực tế bằng số tiền dự kiến của kỳ & Đã trả & Kỳ trả nợ hoàn tất nghĩa vụ, không còn tính vào dư nợ đến hạn. \\ \hline
\end{longtable}
\end{smalltable}

\subsection{Vòng đời Sổ tiết kiệm (\texttt{Saving})}

\begin{smalltable}
\begin{longtable}{|L{0.15\linewidth}|L{0.19\linewidth}|L{0.19\linewidth}|L{0.13\linewidth}|L{0.27\linewidth}|}
\lthead{\textbf{Trạng thái hiện tại} & \textbf{Sự kiện kích hoạt} & \textbf{Điều kiện chuyển trạng thái} & \textbf{Trạng thái tiếp theo} & \textbf{Ý nghĩa nghiệp vụ}}
Đang gửi & Đến ngày kết thúc kỳ hạn gửi & Ngày hiện tại bằng hoặc lớn hơn ngày đáo hạn & Đã đáo hạn & Sổ tiết kiệm sẵn sàng để rút toàn bộ gốc và lãi hoặc tái tục sang kỳ hạn mới. \\ \hline
\end{longtable}
\end{smalltable}

\section{Bảng đối chiếu ca sử dụng với đối tượng}

Bảng ma trận đối chiếu nhằm kiểm tra tính đầy đủ của mô hình phân tích, bảo đảm mọi ca sử dụng đều có đối tượng tương ứng phụ trách và không có đối tượng nào bị dư thừa trong phạm vi hệ thống:

\begin{smalltable}
\begin{longtable}{|L{0.11\linewidth}|L{0.17\linewidth}|L{0.24\linewidth}|L{0.41\linewidth}|}
\lthead{\textbf{Mã Ca sử dụng} & \textbf{Tên ca sử dụng} & \textbf{Các đối tượng nghiệp vụ tham gia} & \textbf{Vai trò nghiệp vụ của các đối tượng}}
\code{WAL-01} & Xem số dư các ví & \code{Wallet}, \code{Transaction}, \code{Transfer}, \code{DebtPayment} & \code{Wallet} tổng hợp dòng tiền từ \code{Transaction}, luân chuyển từ \code{Transfer} và biến động từ \code{DebtPayment} để hiển thị số dư và tỷ trọng. \\ \hline
\code{TXN-01} & Ghi chép giao dịch thu/chi & \code{Wallet}, \code{Transaction}, \code{Group}, \code{Category} & \code{Transaction} khởi tạo bản ghi thu chi, \code{Category} và \code{Group} cung cấp cấu trúc phân loại, \code{Wallet} điều chỉnh số dư khả dụng. \\ \hline
\code{TXN-02} & Sửa / Xóa giao dịch & \code{Wallet}, \code{Transaction}, \code{Group}, \code{Category} & \code{Transaction} cập nhật hoặc đánh dấu hủy, \code{Wallet} tính toán bù trừ hoàn lại số dư tương ứng. \\ \hline
\code{TXN-03} & Xem lịch sử thu/chi & \code{Wallet}, \code{Transaction}, \code{Group}, \code{Category} & Truy xuất tập hợp các \code{Transaction} theo bộ lọc thời gian, phân loại theo \code{Group} và \code{Category}. \\ \hline
\code{TRF-01} & Chuyển tiền nội bộ giữa ví & \code{Wallet}, \code{Transfer}, \code{Transaction} & \code{Transfer} điều phối luân chuyển tiền giữa hai \code{Wallet}, kích hoạt tạo \code{Transaction} phụ nếu phát sinh phí dịch vụ. \\ \hline
\code{CAT-01} & Phân loại thu/chi theo danh mục & \code{Transaction}, \code{Group}, \code{Category} & Ánh xạ sự kiện chi tiêu vào cơ cấu danh mục phân tầng để phục vụ theo dõi ngân sách. \\ \hline
\code{CAT-02} & Quản lý danh mục phân loại & \code{Group}, \code{Category} & Quản trị dữ liệu danh mục gồm thêm mới, cập nhật thông tin hiển thị và vô hiệu hóa danh mục cũ. \\ \hline
\code{DSH-01} & Xem dashboard tổng quan & \code{Wallet}, \code{Transaction}, \code{Group}, \code{Category}, \code{DebtLedger}, \code{GoalFund} & Tổng hợp số liệu tài sản từ \code{Wallet}, luồng tiền từ \code{Transaction}, nghĩa vụ ròng từ \code{DebtLedger} và tiến độ từ \code{GoalFund}. \\ \hline
\code{DBT-01} & Quản lý công nợ & \code{Debt}, \code{RepaymentInstallment}, \code{DebtPayment}, \code{Partner}, \code{Wallet} & \code{Debt} quản lý thỏa thuận nợ, \code{RepaymentInstallment} chia kỳ, \code{DebtPayment} ghi nhận thanh toán và tác động đến \code{Wallet}. \\ \hline
\code{FR-DBT-03} & Sao kê công nợ theo người & \code{Partner}, \code{Debt}, \code{DebtLedger} & \code{Partner} cung cấp thông tin đối tác, \code{DebtLedger} tổng hợp toàn bộ các \code{Debt} để kết xuất sao kê phải thu và phải trả. \\ \hline
\code{SAV-01} & Quản lý sổ tiết kiệm & \code{Saving}, \code{GoalFund}, \code{Wallet} & \code{Saving} quản lý thông tin tiền gửi và tính lãi, \code{GoalFund} tích hợp giá trị tài sản, \code{Wallet} là nguồn trích tiền gửi ban đầu. \\ \hline
\code{FCT-01} & Quản lý chứng chỉ quỹ đầu tư & \code{FundCertificate}, \code{GoalFund}, \code{Wallet} & \code{FundCertificate} theo dõi danh mục đầu tư và lãi lỗ thị trường, \code{GoalFund} tích hợp vào quỹ mục tiêu, \code{Wallet} thanh toán mua. \\ \hline
\code{GFD-01} & Phân bổ \& theo dõi quỹ mục tiêu & \code{GoalFund}, \code{Saving}, \code{FundCertificate}, \code{Wallet} & \code{GoalFund} nhận phân bổ tiền mặt từ \code{Wallet}, tổng hợp giá trị từ \code{Saving} và \code{FundCertificate} để tính toán tiến độ hoàn thành. \\ \hline
\end{longtable}
\end{smalltable}

\section{Bảng tự kiểm tra chất lượng phân tích}

\subsection{Bảng tự kiểm tra chất lượng phân tích Milestone 2}

\begin{smalltable}
\begin{longtable}{|L{0.29\linewidth}|L{0.24\linewidth}|L{0.09\linewidth}|L{0.31\linewidth}|}
\lthead{\textbf{Câu hỏi tự kiểm tra} & \textbf{Tiêu chí đạt yêu cầu} & \textbf{Đánh giá} & \textbf{Ghi chú minh chứng}}
Báo cáo có bị lẫn các lớp kỹ thuật (Controller, Service, Repository, DAO) không? & Tuyệt đối không xuất hiện các lớp kỹ thuật này & Đạt & Toàn bộ 13 đối tượng đều là thực thể đời thực hoặc đối tượng giá trị miền tài chính. \\ \hline
Báo cáo có bị lẫn chi tiết giao diện người dùng không? & Không đề cập nút bấm, form nhập liệu hay cửa sổ & Đạt & Các danh từ giao diện (form, button) đã được lọc bỏ hoàn toàn tại mục 2.2. \\ \hline
Tên đối tượng có bám sát ngữ cảnh đời thực và Use Case không? & Bám sát thuật ngữ trong tài liệu req-uc.md và file Excel & Đạt & 13 đối tượng phản ánh chính xác các khái niệm tiền tệ, sổ sách, công nợ và đầu tư. \\ \hline
Thẻ CRC có phân định rõ ràng việc đối tượng biết và việc đối tượng tự làm không? & Có sự tách bạch rành mạch giữa Knowing và Doing & Đạt & Cột Responsibilities trong 13 thẻ CRC tại mục 5 phân định rõ trách nhiệm thông tin và hành vi. \\ \hline
Các phép tính tiền tệ có quy định sử dụng số thập phân chính xác không? & Bắt buộc quy ước kiểu dữ liệu BigDecimal & Đạt & Quy tắc \code{BR-SYS-01} và bảng từ điển mục 3 quy định rõ sử dụng \code{BigDecimal}. \\ \hline
Phạm vi dữ liệu giữa các phân hệ có bị lấn sân nhau không? & Mỗi đối tượng chỉ quản lý dữ liệu tự thân, điểm chạm ở mức khái quát & Đạt & Mục 1.3 và 4.2 xác lập ranh giới điểm chạm rõ ràng qua giao dịch tiền tệ chung. \\ \hline
Có bị lạm dụng các sơ đồ UML chi tiết của pha thiết kế không? & Không sử dụng sơ đồ lớp chi tiết có kiểu dữ liệu lập trình hay sơ đồ tuần tự hàm & Đạt & Sử dụng mô hình trực quan 2 tầng dạng khối khái niệm Mermaid đơn giản. \\ \hline
\end{longtable}
\end{smalltable}

\subsection{Đối chiếu tuân thủ 10 Quality Gates kỹ thuật}

\begin{smalltable}
\begin{longtable}{|L{0.20\linewidth}|L{0.26\linewidth}|L{0.09\linewidth}|L{0.38\linewidth}|}
\lthead{\textbf{Cổng kiểm định} & \textbf{Tiêu chí kỹ thuật} & \textbf{Kết quả} & \textbf{Ghi chú xác thực}}
Gate 1: Ranh giới miền bài toán & Không rò rỉ thuật ngữ thiết kế OOD vào tài liệu OOA & Đạt & Không có Controller, Service, Repository, DAO, DTO hay cấu trúc bảng SQL. \\ \hline
Gate 2: Nhất quán góc nhìn Diataxis & Tài liệu thuần túy thuộc góc nhìn Diataxis: Tài liệu giải thích và đặc tả miền (Explanation \& Reference) & Đạt & Bố cục chuẩn xác theo 9 mục quy định tại framework.md, không giảng bài lý thuyết suông. \\ \hline
Gate 3: Độ thuần nhất ngôn ngữ & Thuật ngữ kỹ thuật và văn bản tiếng Việt tự nhiên, không chèn dịch nghĩa kép vụng về & Đạt & Sử dụng thuật ngữ chuyên môn kế toán và tài chính chuẩn mực. \\ \hline
Gate 4: Bộ lọc văn phong máy học & Loại bỏ hoàn toàn gạch ngang em-dash, mẫu câu đối lập giả tạo và bôi đậm đầu dòng cơ học & Đạt & Số lượng em-dash bằng 0; không có mẫu câu ``không phải X mà là Y''; cấu trúc danh sách chuyển thành bảng biểu. \\ \hline
Gate 5: Định chuẩn định danh mục & Đánh số thứ tự và mã định danh chuẩn mực & Đạt & Các quy tắc nghiệp vụ mang mã chuẩn \code{BR-SYS}, \code{BR-WAL}, \code{BR-TXN}, \code{BR-TRF}, \code{BR-CAT}, \code{BR-DBT}, \code{BR-PRT}, \code{BR-SAV}, \code{BR-FCT}, \code{BR-GFD}. \\ \hline
Gate 6: Tính toàn vẹn công thức toán & Không chứa ký tự sổ dọc trần trong các công thức toán học & Đạt & Các biểu thức công thức sử dụng ký tự toán học chuẩn, không xung đột bảng Markdown. \\ \hline
Gate 7: Vệ sinh khối mã và ký tự & Khối mã Mermaid sạch, không có khoảng trắng sau dấu backtick & Đạt & Toàn bộ các khối mermaid được định dạng chuẩn xác. \\ \hline
Gate 8: Mô hình hóa trực quan hai tầng & Mỗi sơ đồ không vượt quá giới hạn nhận thức 9 node & Đạt & Sơ đồ Tầng 1 (5 nodes), Tầng 2a (5 nodes), Tầng 2b (8 nodes), 100\% tuân thủ giới hạn nhỏ hơn hoặc bằng 9 nodes. \\ \hline
Gate 9: Ma trận truy vết khép kín & 100\% ca sử dụng được ánh xạ đầy đủ với các đối tượng miền & Đạt & Ma trận mục 8 ánh xạ trọn vẹn 11 ca sử dụng chính và 1 yêu cầu chức năng sao kê. \\ \hline
Gate 10: Ranh giới thực tế và đánh đổi & Xác định rõ giới hạn của mô hình phân tích miền & Đạt & Ranh giới phân tích được xác lập rõ ràng tại mục 1.3, làm tiền đề chuyển giao cho pha OOD. \\ \hline
\end{longtable}
\end{smalltable}

\end{document}